\documentclass[11pt,a4paper]{article}
\usepackage[utf8]{inputenc}
\usepackage[T1]{fontenc}
\usepackage[spanish]{babel}
\usepackage{lmodern}
\usepackage{microtype}
\usepackage{enumitem}
\usepackage{hyperref}
\hypersetup{colorlinks=true,linkcolor=blue,urlcolor=blue}
\usepackage{parskip}

\title{Documento general de la tesis\\[2pt]
\large Sintesis dimensional cinematica de un exoesqueleto de rehabilitacion de dedo}
\author{}
\date{}

\begin{document}
\maketitle

Este documento resume de que trata la tesis, cual es el objetivo, en que se basa
y que falta. Acompana a la entrega de codigo organizada en tres carpetas
descargables dentro de esta misma ubicacion (\texttt{entrega\_codigo/}).

\section{De que trata la tesis}

La tesis aborda la \textbf{sintesis dimensional cinematica} de un exoesqueleto de
rehabilitacion para un dedo de la mano. El objetivo del diseno es que el
exoesqueleto reproduzca el \textbf{agarre de pinza fina} que se midio por captura
de movimiento (mocap) en un dedo humano real. Es decir, dado un patron de
movimiento humano (la trayectoria de las falanges durante la pinza fina), se
buscan las dimensiones de los eslabones del mecanismo del exoesqueleto para que
sus falanges sigan ese mismo movimiento.

El problema se plantea como una optimizacion: se define una funcion objetivo que
mide cuanto se parece la trayectoria del exoesqueleto a la del mocap, y se buscan
los parametros dimensionales que la minimizan.

Se trabajan dos versiones del modelo:

\begin{itemize}[leftmargin=1.4em]
  \item \textbf{Modelo simplificado (parcial), 16 parametros:} cubre la
        cinematica hasta la falange medial (articulaciones MCP e IFP). Es el
        banco de pruebas para comparar metaheuristicas de forma justa.
  \item \textbf{Modelo completo, 21 parametros:} cubre las tres falanges
        (proximal, medial y distal) con cinco mecanismos. Extiende el modelo
        simplificado para incluir la falange distal.
\end{itemize}

\section{Objetivo}

El objetivo metodologico es \textbf{comparar dos metaheuristicas de forma justa y
simetrica}:

\begin{itemize}[leftmargin=1.4em]
  \item Evolucion Diferencial (ED)
  \item Algoritmo Genetico (GA)
\end{itemize}

La comparacion es justa porque ambos algoritmos:

\begin{itemize}[leftmargin=1.4em]
  \item Minimizan \textbf{la misma funcion objetivo} con los mismos pesos.
  \item Usan \textbf{los mismos limites} (bounds) de los parametros.
  \item \textbf{Sintonizan sus hiperparametros con Optuna} (muestreador TPE, 25
        trials, semilla 42), de modo que ninguno queda en desventaja por una mala
        eleccion manual de hiperparametros.
\end{itemize}

Asi, la diferencia de resultados mide el ALGORITMO y no la formulacion del
problema. El ganador de la comparacion en el modelo simplificado (ED, con error
global 2.17~mm frente a 2.55~mm del GA) es la metaheuristica que despues se usa
para optimizar el \textbf{modelo completo}.

\section{En que se basa}

\begin{itemize}[leftmargin=1.4em]
  \item \textbf{Datos de mocap de pinza fina (120 puntos):} la trayectoria
        objetivo que el exoesqueleto debe reproducir
        (\texttt{mocap\_pinza\_fina\_120pts.csv}).
  \item \textbf{Distancia de Chamfer} como metrica de forma: compara la nube de
        puntos de la trayectoria del exoesqueleto con la del mocap, ponderada por
        articulacion.
  \item \textbf{Alineacion rigida optima tipo Kabsch/Procrustes:} antes de medir
        el error se alinea rigidamente la trayectoria simulada con la del mocap,
        de modo que se compara la FORMA del movimiento y no su posicion absoluta.
  \item \textbf{Penalizacion de monotonicidad (W\_MONO = 5.0):} obliga al
        movimiento a avanzar de forma coherente, sin retrocesos artificiales.
  \item \textbf{Marco de sintesis de mecanismos de 5 y 4 barras:} las etapas del
        exoesqueleto se modelan como combinaciones de mecanismos de cinco y
        cuatro barras, resueltos con los solucionadores de \texttt{comun.py}
        (\texttt{sol\_5\_barras}, \texttt{solve\_four\_bar},
        \texttt{circle\_intersections}).
  \item \textbf{Restricciones cinematicas + penalizacion fija de 1000:} los
        disenos que no ensamblan o violan restricciones fisiologicas reciben una
        penalizacion alta.
  \item \textbf{El articulo COMROB} como soporte y marco de publicacion de los
        resultados.
\end{itemize}

Detalle de la comparacion completa vs parcial (funcion objetivo, monotonicidad,
Optuna y restricciones) en el reporte
\texttt{AUDITORIA\_completa\_vs\_parcial.tex}.

\section{Tercera etapa para la falange distal}

Al modelo original se le \textbf{agrego una TERCERA ETAPA de mecanismo de cuatro
barras} para dar movimiento a la \textbf{falange distal respecto de la proximal},
y esa tercera etapa se \textbf{incorporo al analisis cinematico} (queda dentro
del barrido y de la funcion objetivo del modelo completo).

\paragraph{Aclaracion anatomica (no proviene del usuario).} En la anatomia del
dedo, la falange distal articula directamente con la falange media (o medial) a
traves de la articulacion interfalangica distal (DIP/IFD). Se deja constancia de
esta precision como aclaracion separada, sin sustituir la frase del usuario
(``respecto a la proximal''). Confirmar con el usuario la referencia anatomica que
debe quedar en la redaccion final de la tesis.

En el codigo (\texttt{modelo\_completo.py}, seccion ``TERCER MECANISMO DE 4
BARRAS'') esta tercera etapa se implementa como un mecanismo de \textbf{cuatro
barras}:

\begin{itemize}[leftmargin=1.4em]
  \item Acoplador \texttt{Link9\_3}.
  \item Balancin \texttt{Link10\_3}.
  \item Soporte dorsal definido por \texttt{back3\_3} (retroceso a lo largo del
        eje de la falange) y \texttt{up3\_3} (standoff dorsal), que fija el punto
        de anclaje.
  \item El punto \texttt{D3} se obtiene por \textbf{interseccion de circulos}
        (\texttt{circle\_intersections}), y de ahi se calcula el angulo de la
        falange distal.
\end{itemize}

Incluye ademas el tope fisiologico \texttt{DIP\_MAX\_DEG = 35} grados y un termino
de perfil angular DIP en la funcion objetivo.

\section{Que falta}

\begin{itemize}[leftmargin=1.4em]
  \item \textbf{Construir el CAD} del exoesqueleto con los 21 parametros
        optimizados del modelo completo.
  \item \textbf{Verificar el agarre} del CAD contra las trayectorias alineadas
        del mocap.
  \item \textbf{Ajustar por interferencias mecanicas} (colisiones y choques entre
        eslabones que la cinematica idealizada no captura).
\end{itemize}

\section{Estructura de la entrega}

\begin{itemize}[leftmargin=1.4em]
  \item \texttt{01\_optimizacion\_parcial\_es/}: optimizacion parcial (modelo
        simplificado, 16 parametros) en espanol, con comparacion ED vs GA,
        figuras y resultados.
  \item \texttt{02\_optimizacion\_completa\_es/}: optimizacion completa (modelo de
        3 falanges, 21 parametros) en espanol, con la via canonica
        (\texttt{optimizar\_completo\_ED.py}) y la via legacy de referencia del
        tercer mecanismo.
  \item \texttt{03\_optimizacion\_parcial\_en/}: version en ingles de la
        optimizacion parcial.
  \item \texttt{AUDITORIA\_completa\_vs\_parcial.tex}: reporte de auditoria que
        verifica que la optimizacion completa incorpora las mismas mejoras que la
        parcial.
\end{itemize}

Cada carpeta es autocontenida (incluye su CSV de mocap y sus scripts) y trae su
propio documento de texto (\texttt{LEEME.txt} / \texttt{README.txt}) con el
detalle de scripts, dependencias exactas, orden de ejecucion y resultados
esperados.

\end{document}
